\chapter{Resultados}
\label{cap:resultados}

% =====================================================================================
% CAPITULO BLOQUEADO. NO ESCRIBIR TODAVIA.
%
% Requiere: (1) las 369 muestras de 15 qubits, generandose en la nube;
%           (2) los tres modelos entrenados;
%           (3) las cifras de la EDA rehechas sobre el dataset v2.
%
% Ver doc/memoria/pendientes.md.
% =====================================================================================
%
% ⚠️ REGLA DE ESTE CAPITULO, tomada de las Instrucciones seccion 2.6:
%    «Es una exposicion OBJETIVA, sin valorar los resultados ni justificarlos».
%    Aqui van tablas, graficas y datos. Toda interpretacion —por que un modelo gana,
%    que significa, ventajas y desventajas— va al capitulo 11.
%    Si al redactar aparece un «porque» o un «esto sugiere», va en el 11, no aqui.

\section{Comparativa global}
\label{sec:comparativa}
Los tres modelos sobre el mismo protocolo. Tabla principal de la memoria.

\section{Desglose por eje OOD}
\label{sec:desglose}
% Un parrafo introductorio antes de las subsecciones: la norma prohibe dos encabezados
% consecutivos sin texto entre medias.
\subsection{Eje de tipo: generalización a algoritmos nunca vistos}
\subsection{Eje de tamaño: extrapolación}
\subsection{Eje temporal: robustez a la deriva del hardware}

\section{Desglose por observable}
\label{sec:por_observable}
Qué observables se predicen bien y cuáles no. Reportar el $n$ de cada celda.

\section{Comportamiento frente al listón}
\label{sec:vs_baseline}
% Mejora relativa frente a NO MITIGAR, que es el criterio de exito declarado en el
% objetivo general del capitulo 5. Cifras, sin valoracion.

\section{Distribución del error}
\label{sec:errores}
Dónde falla cada modelo y si los tres fallan en las mismas muestras. Medianas y
percentiles además de las medias, por las colas pesadas.
